\noindent
A companion paper described gradient descent on Boolean targets as noise-limited greedy growth of features. Here we ask whether a \emph{discrete} learner, one with no continuous feature parameters, can do at least as well as gradient-descent-trained networks at equal compute. We ask it first on the Boolean toy targets of that study and then on standard benchmarks. There are four steps.

\emph{A discrete learner for the toy targets.} DFG (discrete feature growth) keeps the networks' strategy and replaces their search. It grows Fourier characters greedily, testing each against the noise of the current residual, and it makes compositions of features already found cheap. It also fits a linear background and reads out with the sign of a fit. Its search is exhaustive enumeration with union-bound tests, run within a fixed share of one network run's FLOPs. The references are $\mu$P two-hidden-layer MLPs trained by Adam on the same training sets, at $d=30$. On a blind suite of 7 held-out targets, DFG reaches 10\% test error with less data on all 7: a median factor of 3.1, with at most 0.26 of one network run's FLOPs per fit. On 41 development targets the median factor is 5--6, and it grows with the input dimension. Heuristically, enumeration pays for search in compute and about $2k\ln d$ samples for a $k$-sparse parity, where the networks pay about $d^{k-1}$ samples.

\emph{The same learner on standard tasks.} Three independently built adaptations of DFG were run on 18 classification tasks (MNIST, Fashion-MNIST and 16 PMLB datasets). On tabular data they are about level with a plain MLP, and their error is 12--17\% higher than an L2-regularized MLP's at $n\le3000$. On images they are clearly behind at large $n$, closer to logistic regression than to the MLP.

\emph{A discrete portfolio for both.} A pre-registered portfolio stacks DFG, LightGBM, ExtraTrees, nearest-neighbour votes and logistic regression on out-of-fold predictions. It was compared with the networks and MLPs on two fresh sets fixed before any design.
\begin{itemize}
\item \emph{Fresh toy suite.} It needs clearly less data than the per-target best of three network configurations on 9 of 10 targets. The mean $\log_2$ ratio of the data needed is $-1.41$. It needs more data on one symmetric function, $[|x_1+\dots+x_9|\le1]$ (1229 samples against 350).
\item \emph{Fresh tabular datasets.} Its mean error is below the hindsight-best MLP's on 8 of 10. This is not significant (sign test $p=0.11$).
\item \emph{Images.} It is slightly ahead of every MLP up to $n=3000$, with 1 of 8 cells resolved. At $n\ge30000$ it is clearly behind a regularized ensemble of three MLPs.
\end{itemize}
Post-hoc ablations put the toy margins on the character search and the real-data margins on the classical stack: the portfolio is an AutoML-style dispatch by input type, not one discrete mechanism.

\emph{Sum features and permutation-invariant image members.} Part 9 targets the two losses. It adds features that are functions of signed sums of inputs, admitted by a union bound over all subsets and signs, and an image member built from subset sums over pixel sets chosen by correlation, so that it never uses the pixel layout. On a held-out suite of 12 Boolean targets whose family was written by an agent blind to the design, the new portfolio needs clearly less data than the best of the pre-registered network variants on 10 targets and more on none (mean $\log_2$ ratio $-1.45$; $-1.25$, with 2 targets worse, when two more weight decays are added to the networks). On the 9 sum-structured targets Part 8's portfolio is behind, at $+0.63$ against the new portfolio's $-1.53$ ($+0.91$ against $-1.26$ with the two added variants). On two fresh 10-class QuickDraw image sets it beats a strong MLP ensemble at $n\ge10000$ (mean $\log_2$ error ratio $-0.18$) and has the lowest error of every permutation-invariant reference at 13 of the 14 cells (7 training sizes on each set), but a CPU-matched ensemble of wider MLPs ties it at $n=60000$, and on one set has slightly lower error there (14.8\% against 14.9\%). Controls qualify the image result. The pooling sets recover the pixel grid from correlations; an MLP given the same pooled sums matches the portfolio at $n=60000$; the winning member is gradient-boosted trees, whose gain splits between the pooled features and longer boosting; a fit at $n=60000$ takes 20--30 CPU minutes; and a small CNN that sees the pixel layout, run on one set at two sizes, has lower error at both, clearly so at $n=10000$.

Independent audits checked leakage, compute accounting, discreteness and the strength of the references in Parts 8 and 9. Their findings strengthened the references, added a blind test suite and post-hoc controls, and changed the wording of the claims; no learner was changed after its freeze except for one recorded crash fix.
